% Generated draft: requires mathematical and source review.
\documentclass[11pt]{article}
\usepackage[margin=25mm]{geometry}
\usepackage{fontspec}
\setmainfont{Latin Modern Roman}
\usepackage{amsmath,amssymb,mathtools,unicode-math}
\setmathfont{Latin Modern Math}
\usepackage{enumitem,needspace,xurl,hyperref}
\setlength{\parindent}{0pt}
\setlength{\parskip}{5pt}
\setlength{\emergencystretch}{3em}
\newcommand{\N}{\mathbb{N}}
\newcommand{\Z}{\mathbb{Z}}
\newcommand{\Q}{\mathbb{Q}}
\newcommand{\R}{\mathbb{R}}
\newcommand{\C}{\mathbb{C}}
\newcommand{\F}{\mathbb{F}}
\newcommand{\A}{\mathbb{A}}
\newcommand{\PP}{\mathbb{P}}
\newcommand{\E}{\mathbb{E}}
\newcommand{\eps}{\varepsilon}
\newcommand{\dd}{\,\mathrm{d}}
\begin{document}
\section*{UnsolvedMath ID: 30005067}
\textbf{Source problem code:} catalogue entry
\subsection{Definitions and mathematical statement}
The supplied record asks for a classification of radial solutions of
\[
(-\Delta)^\gamma u=A u^p\qquad\text{in }\R^n\setminus\{0\},
\]
with an isolated, non-removable singularity at the origin. Here radial means that \(u(x)=U(|x|)\) for a profile \(U\). The notation \(( -\Delta)^\gamma\) denotes the fractional Laplacian; on functions for which the Fourier definition applies, its symbol is \(|\xi|^{2\gamma}\).

To state the classification target without supplying missing hypotheses, let \(\mathcal X\) denote the intended admissible solution class, including the interpretation of the fractional Laplacian, the meaning of the equation, and any sign, regularity, integrability, and asymptotic conditions. For each intended admissible parameter tuple \((n,\gamma,p,A)\), define
\[
\mathcal S_{n,\gamma,p,A}(\mathcal X)
=\left\{u\in\mathcal X:\begin{array}{l}
u(x)=U(|x|),\\
(-\Delta)^\gamma u=A u^p\ \text{on }\R^n\setminus\{0\},\\
0\text{ is an isolated, non-removable singularity of }u
\end{array}\right\}.
\]
Non-removability is understood relative to the chosen solution class: the solution does not extend across the origin as a solution in that class. The problem is to describe every member of this set for every parameter tuple allowed by the source. The record does not determine \(\mathcal X\) or the complete parameter range, so this is a conditional classification target requiring clarification before it becomes a fully specified mathematical problem.

The source requires \(A\) to be normalized so that a designated power-law function \(u_0\) is a singular solution; \(A\) is therefore not presented as an arbitrary independent constant. The displayed power and parameter intervals are corrupted in the supplied original extract. For \(p\ne1\), the scaling-compatible reconstruction is
\[
\alpha=\frac{2\gamma}{p-1},\qquad u_0(r)=r^{-\alpha}.
\]
This reconstruction and its admissibility require source verification. Under this reconstruction, cylindrical coordinates may be defined by
\[
t=-\log r,\qquad v(t)=e^{-\alpha t}U(e^{-t}),\qquad t\in\R.
\]
The supplied clean statement describes the resulting equation as autonomous and nonlocal, and presents its classification as equivalent to the radial classification. The referenced equation (3), its operator, and its admissible solution space are not supplied, so no explicit version of that equation is asserted here.

The extract additionally discusses subcritical radial solutions with fast decay, their existence as building blocks, and mapping properties of the linearized operator. It also asserts that non-removable singularities have the growth rate of the designated power solution. The hypotheses and exact meanings of fast decay and exact growth rate are not recoverable from this record; these assertions are retained as qualified source context, not added as established properties of every member of \(\mathcal S_{n,\gamma,p,A}(\mathcal X)\).
\subsection{Short English statement}
Classify radial solutions of \(( -\Delta)^\gamma u=A u^p\) on \(\R^n\setminus\{0\}\) with an isolated, non-removable singularity at the origin, using the source's prescribed power-solution normalization of \(A\). The supplied record also describes an equivalent autonomous nonlocal equation in cylindrical coordinates. Corrupted parameter ranges, unspecified solution and decay conditions, and the missing cylindrical equation prevent a fully determined formulation. The supplied status is partially solved; this is dated curation requiring independent review.
\subsection{Sources}
\begin{itemize}
\item Supplied local UnsolvedMath registry record, permanent ID \(30005067\), title: Classification of Singular Radial Fractional-Laplacian Solutions. Exact source problem code: \texttt{catalogue entry}. No individual proposer is identified in the record; the supplied proposed year is 2022.
\item Supplied report citation: MFO-RIMS Tandem Workshop: Nonlocality in Analysis, Probability and Statistics (2023), Oberwolfach Reports, DOI 10.4171/owr/2022/15, pp.\ 709--715. \url{https://doi.org/10.4171/owr/2022/15}. The record also labels this reference Oberwolfach Report 15/2022. The discrepancy between the citation's 2023 designation and its 2022 identifiers is unresolved.
\item Supplied literature source: Ao--Chan--DelaTorre--Fontelos--Gonzalez--Wei, ODE methods. \url{https://arxiv.org/abs/1910.14512}. The record supplies no fuller bibliographic entry.
\item Literature assessment date: \texttt{2026-08-22}. This is supplied dated curation, not an independently verified assessment. The supplied status is partially solved. The supplied assessment says that the cited primary work develops an infinite-ODE method and proves nondegeneracy and related structural results for fast-decay singular solutions, while the full classification of the nonlocal ODE remains incomplete. The background expressly requests independent verification of time-sensitive status and cited claims.
\item The registry labels the statement \texttt{corrected\_verified} and reports that it was reconstructed and checked against the official Oberwolfach report. These are supplied provenance claims; no source visit or independent verification is claimed in this draft.
\end{itemize}
\end{document}
